\section{从“万物哈密顿量”到常用模型}

多体理论的出发点往往是电子与离子的完整哈密顿量。通过冻结离子、保留振动、投影到 Bloch/Wannier 基等步骤，可得到 jellium、电子–声子、紧束缚、Hubbard、Anderson/Kondo 等模型。层次关系如下（对齐施均仁讲义第 1.4.2 节）。

\subsection{层次概览}

\begin{center}
\begin{tabular}{c}
电子+离子的“万物”哈密顿量\\
$\downarrow$ 冻结离子 / $\downarrow$ 保留晶格振动\\
jellium 电子液 \qquad / \qquad 电子–声子耦合\\
$\downarrow$ 周期势 $\to$ Bloch / Wannier\\
紧束缚模型\\
$\downarrow$ 单带、局域相互作用\\
Hubbard 模型\\
$\downarrow$ 杂质极限\\
Anderson 杂质 $\to$ Kondo 模型
\end{tabular}
\end{center}

\subsection{紧束缚与 Hubbard}

Wannier 基下单带紧束缚
\begin{equation}
  H=-\sum_{ij\sigma}t_{ij}c_{i\sigma}^\dagger c_{j\sigma}
  +\sum_{ii'jj',\sigma\sigma'}
  U_{ii'jj'}
  c_{i\sigma}^\dagger c_{i'\sigma'}^\dagger
  c_{j'\sigma'}c_{j\sigma}.
\end{equation}
若只保留近邻跃迁与同址排斥，即得 Hubbard 模型
\begin{equation}
  H=-t\sum_{\langle ij\rangle\sigma}
  \bigl(c_{i\sigma}^\dagger c_{j\sigma}+\mathrm{h.c.}\bigr)
  +U\sum_i n_{i\uparrow}n_{i\downarrow}.
\label{eq:hubbard}
\end{equation}
交换型相互作用可写为自旋形式
\begin{equation}
  H_J=-2\sum_{ij}J_{ij}
  \Bigl(
  \mathbf{S}_i\cdot\mathbf{S}_j+\tfrac14 n_i n_j
  \Bigr),
  \qquad
  \mathbf{S}_i=\tfrac12\sum_{\sigma\sigma'}
  c_{i\sigma}^\dagger\boldsymbol{\tau}_{\sigma\sigma'}c_{i\sigma'}.
\end{equation}

\subsection{Anderson 杂质与 Kondo}

杂质能级与传导带杂化：
\begin{equation}
\begin{aligned}
  H
  &=\sum_\sigma\varepsilon_f f_\sigma^\dagger f_\sigma
  +U n_{f\uparrow}n_{f\downarrow}
  +\sum_{\mathbf{k}\sigma}
  \bigl(V_{\mathbf{k}}f_\sigma^\dagger c_{\mathbf{k}\sigma}+\mathrm{h.c.}\bigr)\\
  &\quad+\sum_{\mathbf{k}\sigma}\varepsilon_{\mathbf{k}}
  c_{\mathbf{k}\sigma}^\dagger c_{\mathbf{k}\sigma}.
\end{aligned}
\end{equation}
深能级、大 $U$ 极限下可积出电荷涨落，得到 Kondo 交换
\begin{equation}
  H=\sum_{\mathbf{k}\sigma}\varepsilon_{\mathbf{k}}
  c_{\mathbf{k}\sigma}^\dagger c_{\mathbf{k}\sigma}
  -J\,\mathbf{S}_{\mathrm{imp}}\cdot\mathbf{s}(0).
\end{equation}
这是重费米子与局域关联效应的重要微观起源。

\begin{note}
Bloch 基下动量守恒可差一个倒格矢 $\mathbf{K}$，对应 Umklapp 过程，对电阻与一维/准一维系统尤其重要。
\end{note}
